% Folder 136, batch 19 — archivists' pages 361–375 (PDF pages 362–376).
% Pass: Opus 5.5 (claude-opus-5-5), 2026-10-03 — first pass, unchecked against the pages by a human.
%
% \page{N} follows the batch numbering (page 1 of the batch = 361). The pencil
% numbers on the leaves run two or three ahead (372 on page 370, 376 on page
% 373, 378 on page 375); the batch numbering, not the pencil, is used here.
% Page 368 is a blank music-ruled sheet (only a few strokes on the edge of an
% underlying leaf show) and is absent.
\documentclass[11pt,a4paper]{article}
\input{../preamble/grothendieck}

\folder{136}
\batch{19}
\pages{361}{375}
\foldertitle{Complexe de De Rham à puissance divisée [conférence de 1976 à l’IHÉS] : notes manuscrites (s.d.).}
\dating{[à partir de 1975-1976]}
\watermark{Édition de démonstration}

\begin{document}

\section{Formule de décalage (Quillen--Bousfield)}

\page{361}
\note{p. 1 de l'auteur ; les feuillets 361 à 365 portent sa numérotation 1 à 5.}

Formule de Decalage (Quillen -- Bousfield) [cf \struck{\ill{}} Thèse Illusie].

$L_{\ast}F(M,n)$ \qquad $L_{\ast}F(M[n])$

\[
  L_i\,\mathrm{Sym}^i(M,n) \simeq L_{i+n}\,\Lambda^i(M,
  \text{\struck{$n+1$}})
\]
\[
  L_i\,\Lambda^i(M,n) \simeq L_{i+n}\,\Gamma^i(M,n+1)
\]
\note{ces deux lignes sont traversées par de longs traits obliques, qui peuvent valoir rature ; la ligne suivante est entièrement biffée : \struck{\ill{}}.}

\[
  M \to M \qquad
  K(M,n) \longrightarrow E(M,n+1) \longrightarrow K(M,n+1)
\]
\note{la flèche $E(M,n+1) \to K(M,n+1)$ est verticale sur la page.}

\[
  \mathrm{Sym}\,K(M,n) \otimes \Lambda E(M,n+1)
\]
\[
  \Lambda^{n+1}\Psi \otimes_{\mathbb{Z}} M
\]
\note{la lecture $\Lambda^{n+1}\Psi$ est \uncertain{douteuse} ; elle est écrite sous $\Lambda E(M,n+1)$.}

\[
  0 \to \mathrm{Sym}^i K(M,n) \to \mathrm{Sym}^{i-1}K(M,n)\otimes\Lambda^1E(M,n+1)
  \to \mathrm{Sym}^{i-2}K(M,n)\otimes\Lambda^2E(M,n+1)
\]
\[
  \to \cdots \to \Lambda^iE(M,n+1) \to \Lambda^iK(M,n+1) \to 0
\]
\[
  0 \to \Lambda^iK(M,n) \to \Lambda^{i-1}K(M,n)\otimes\mathrm{Sym}^1E(M,n+1)
  \to \cdots
\]
\[
  \cdots \to \mathrm{Sym}^iE(M,n+1) \to \mathrm{Sym}^iK(M,n+1) \to 0
\]
\[
  0 \to \Gamma^iK(M,n) \to \Gamma^{i-1}K(M,n)\otimes\Lambda^1E(\ \cdots)
\]

Encadré :
\[
  L_r\,\mathrm{Sym}^i K(M,n+1) \simeq L_{r-i}\,\Lambda^iK(M,n)
\]

\begin{align*}
  L_{\ast}\,\mathrm{Sym}^iX[1] &\simeq L_{\ast}\,\Lambda^iX[-i] \\
  L_{\ast}\,\Lambda^i(X[1]) &\simeq L_{\ast}\,\Gamma^iX[-i] \\
  L_{\ast}\,\mathrm{Sym}^iX[2] &\simeq L_{\ast}\,\Gamma^iX[-2i]
\end{align*}
\note{dans la première ligne, l'indice de $L$ à droite est surchargé (\uncertain{$\ast - i$} corrigé en $\ast$).}

\page{362}
\note{p. 2 de l'auteur.}

(Th) Dold Puppe Ann Inst Fourier.
\[
  L_r\,\mathrm{Sym}^{\ast}_{\uncertain{e}} \simeq L_r\,\mathbb{Z}[\quad]
\]
\note{l'indice de $\mathrm{Sym}^{\ast}$ est \uncertain{un $e$ ou un $\mathbb{Z}$ minuscule}.}

$\Gamma$ \qquad $M_{\bullet}$ un \add{complexe de chaînes de} $A$-modules, du moins pour $M$ concentré en un seul degré.

\struck{$K(M_{\bullet})$}

Def
\[
  L_i\,\mathbb{Z}M_{\ast} = H_i(M_{\ast})
\]
\[
  L_i(\mathrm{Sym})M_{\ast}
\]
\note{entre les deux lignes, un signe d'inclusion vertical, ouvert vers le haut (\uncertain{$\supseteq$}).}

D'où une graduation (dite la graduation par le poids sur l'homologie des esp d'E.M).

Pt clef : Th de Dold-Thom :

$X$ un espace \qquad $\mathrm{SYM}^{\infty}X = {} \cup X \cup X^{(2)} \cup X^{(3)} \cdots$
\[
  X \mapsto \pi_i\,\mathrm{SYM}^{\infty}X \simeq H_iX
  \qquad \mathbb{Z}X \qquad \pi_i\,\mathbb{Z}X
\]
satisfait aux axiomes d'Eilenberg-Steenrod
\note{« axiomes » est écrit sous la ligne, à l'endroit d'un mot biffé \struck{\ill{}}.}

\[
  G = \mathbb{Z}^{k} \times \mathbb{Z}
\]
\note{lecture \uncertain{$\mathbb{Z}^k\times\mathbb{Z}$}, isolée à droite.}

$X$ un espace de Moore $\mathbb{Z}M(G,n)$ \qquad $G$
\[
  K(G,n) \sim \mathrm{Sym}^{\infty}M(G,n) \qquad G \to M(G)
\]
\struck{$M(G),\ M(\ill{})\ \ldots\ \mathrm{Hom}(G,\ill{})$}, \struck{$\mathrm{Ext}(G,\ill{})$}
\note{passage biffé à droite, relié par un trait ; illisible pour l'essentiel.}

\[
  H_i(G,n) = \pi_i\,\mathbb{Z}K(G,n) \overset{\text{Dold-Thom}}{=}
  \pi_i\,\mathbb{Z}\,\mathrm{Sym}^{\infty}M(G,n)
\]
Théorème de Dold. \quad $\mathbb{Z}\,\mathrm{Sym}^{\infty}X \xleftarrow{\ \sim\ } \mathrm{Sym}\,\mathbb{Z}X$
\[
  \pi_i\,\mathrm{Sym}(\mathbb{Z}M(G,n)) = L_i\,\mathrm{Sym}(G)
\]
\note{une flèche verticale à double pointe relie $\pi_i\,\mathbb{Z}\,\mathrm{Sym}^{\infty}M(G,n)$ et $\pi_i\,\mathrm{Sym}(\mathbb{Z}M(G,n))$ ; sous $\mathbb{Z}M(G,n)$ un « $\Vert$ » renvoie à \uncertain{$K(G,n)$}.}

\struck{Free} Par def \quad $\mathbb{Z}M(G,n) \sim K(G,n)$
\[
  H_{n+r}\,K(M,n) \overset{S}{\simeq} H_{n+r+1}\,K(M,n+1)
  \qquad \text{\struck{$n \leq r \leq 2n$}}\ \ r < n.
\]

\section{Opérations cohomologiques stables}

\page{363}
\note{p. 3 de l'auteur.}

\[
  H_{\ast}(K(\mathbb{F}_p,n),\mathbb{F}_p)^{\ast} \simeq
  H^{n+r}(K(\mathbb{F}_p,n),\mathbb{F}_p)
\]
\[
  K(\mathbb{Z},n) \to K(\mathbb{Z},n) \to K(\mathbb{F}_p,n)
\]

Puisque $K(\mathbb{F}_p,n)$ représente la $\mathbb{F}_p$-coh
\note{une flèche descend de $H^{n+r}(K(\mathbb{F}_p,n),\mathbb{F}_p)$ vers « opérations coh de type ».}
opérations coh de type $(\mathbb{F}_p,n\,;\,\mathbb{F}_p,n+r)$, \quad $r<n$.

Opérations stables de deg $r$ \quad $\{P_r\}$ \quad Une collection $\forall n$
\[
  H^n(X,\mathbb{F}_p) \longrightarrow H^{n+r}(X,\mathbb{F}_p)
\]
commute

\begin{tikzcd}
  \widetilde{H}^{n+1}(SX,\mathbb{F}_p) \arrow[r, "P_{n+1}"] &
  \widetilde{H}^{n+1+r}(SX,\mathbb{F}_p) \\
  \widetilde{H}^{n}(X,\mathbb{F}_p) \arrow[r, "P_n"] \arrow[u, "\wr"] &
  \widetilde{H}^{n+r}(X,\mathbb{F}_p) \arrow[u, "\wr"']
\end{tikzcd}

\note{les deux isomorphismes verticaux sont notés « $\wr$ » sans pointe de flèche ; leur sens est ajouté par nous.}

$r<n$
\[
  H^{n+r}(K(\mathbb{F}_p,n),\mathbb{F}_p) \simeq
  H^{n+r+1}(K(\mathbb{F}_p,n+1),\mathbb{F}_p)
\]
\[
  P_n \qquad\qquad P_{n+1}
\]
\[
  H^{n+r-1}(K(\mathbb{F}_p,n-1),\mathbb{F}_p)
\]
\note{ce dernier groupe est écrit à gauche, relié au premier par une flèche.}

\emph{Opérations stables} : Les op stables sont \struck{\ill{}} données par l'algèbre de Steenrod
\[
  p \neq 2 \qquad
  \mathbb{F}_p[P^0, P^1, \ldots, P^n, \ldots, \beta]\,/\,\text{Relations d'Adem}
\]
\[
  P^0 = 1 \qquad \deg P^i = 2i(p-1) \qquad \deg\beta = 1
\]

\marginal{Cohomology operations, Steenrod + Epstein, Ann of Math Studies N° 50, Princeton, p 77}

\page{364}
\note{p. 4 de l'auteur.}

$P^i$ op stables construite par Steenrod.

$p=2$
\[
  \mathbb{F}_2[\mathrm{Sq}^i]\,/\,\mathrm{Sq}^0 = 1 \quad (\mathrm{Sq}^1 = \beta) \quad \text{rel d'Adem}
  \qquad i \geq 0 \qquad \deg \mathrm{Sq}^i = i
\]

\uncertain{$p\ 2$} \qquad
\struck{$\mathrm{Sq}^1\mathrm{Sq}^2 = \mathrm{Sq}$} \quad
$\mathrm{Sq}^1\mathrm{Sq}^2 = \mathrm{Sq}^3 + \mathrm{Sq}^2\mathrm{Sq}^1$ (?) etc …
\note{la relation écrite est fausse telle quelle (Adem donne $\mathrm{Sq}^1\mathrm{Sq}^2 = \mathrm{Sq}^3$) ; le « (?) » est de lui.}

Formule de Cartan :
\[
  P^i(x \times y) = \sum_{j \geq 0} P^j(x)\,P^{i-j}(y)
\]
\[
  P^I = \beta^{\varepsilon_0}P^{\eta_1}\beta^{\varepsilon_1}P^{\eta_2}\cdots
  \qquad I = (\varepsilon_0,\eta_1,\varepsilon_1,\eta_2,\ldots)
  \qquad \varepsilon_i = 0, 1 \qquad \eta_i \geq 0
\]
$\deg I$ de manière évidente ; excès $e(I) = {}$ « $\eta_1 - \sum_{i>1}\eta_i$ ».

admissible
\[
  \eta_i \geq p\,\eta_{i+1} + \varepsilon_i
\]
Les $P^I$ admissibles forment une $\mathbb{F}_p$-base.
\[
  A \longrightarrow H^{\ast}(K(\mathbb{F}_p,n),\mathbb{F}_p)
  \qquad x \mapsto x\,\iota_n
\]

\emph{Théorème} Serre ($p=2$) : $H^{\ast}(K(\mathbb{F}_p,n),\mathbb{F}_p) = {}$ Algèbre anticommutative libre sur $\mathbb{F}_p$ engendrée par les monômes \emph{admissibles} d'excès $<n$.

$H_{\ast}(K(\mathbb{F}_p,n),\mathbb{F}_p) = {}$ Alg anticom à puiss divisée universelle engendrée par les éléments duaux.

\page{365}
\note{p. 5 de l'auteur.}

\emph{Dém 1)} Séminaire Cartan

On regarde la suite spectrale du classifiant
\[
  E_2 = \mathrm{Tor}^{H_{\ast}(K(\mathbb{F}_p,n),\mathbb{F}_p)}(\mathbb{F}_p,\mathbb{F}_p)
  \Longrightarrow H_{\ast}(K(\mathbb{F}_p,n+1),\mathbb{F}_p)
\]
\[
  K(\mathbb{F}_p,n+1) = B_{K(\mathbb{F}_p,n)}
\]
Suite sp d'Eilenberg-Moore, Steenrod - Rothenberg, etc …
\[
  \cdots\ \mathbb{Z}(G\times G) \rightrightarrows \mathbb{Z}G \longrightarrow \mathbb{Z}
  \qquad \mathbb{Z}(G\times G) = \mathbb{Z}G\otimes\mathbb{Z}G
\]
\struck{$\otimes\,\mathbb{Z}K(\mathbb{F}_p,n)$} \quad
$\mathbb{Z}K(\mathbb{F}_p,n) = M$ \quad alg diff graduée augmentée. Bar Constr.
\[
  \cdots\ H_{\ast}(K(\mathbb{F}_p,n))^{\otimes 2} \rightrightarrows
  H_{\ast}K(\mathbb{F}_p,n) \to M
\]
\[
  A = \mathbb{Z}(G) \qquad \mathbb{Z}(EG) = {\ast} \qquad \mathbb{Z}(BG)
\]

2) Méthode de Serre \quad Comm Helv 1954
\[
  H^{\ast}(K(\mathbb{F}_p,n+1),\mathbb{F}_p)\otimes H^{\ast}(K(\mathbb{F}_p,n),\mathbb{F}_p)
  \Longrightarrow H^{\ast}(\ast)
\]

\marginal{Springer : Symposium en l'honneur de Steenrod, Article de May.}

\section{L'algèbre de Steenrod et sa duale}

\page{366}
\note{feuille à portées, écrite en largeur ; la colonne de droite est transcrite d'abord, puis celle de gauche.}

Comult naturelle sur $A$

$P$ opération coh \quad $(\mu^{\ast}P)(xy) = \sum_i P_i x \cdot P'_i y$

Par la formule de Cartan
\[
  \Delta P^i = \sum_{j+k=i} P^j \otimes P^k .
\]
\note{la lecture $\sum_i P_i x\cdot P'_i y$ est \uncertain{incertaine}.}

\emph{Milnor} : \quad $j \geq 1$ \quad $\xi_j$ \quad dual de $P^{p^{j-1}} \cdots P^{p^2}P^{p}P^{p^0}$ \quad (monôme adm)

$j \geq 0$ \quad $\tau_j$ \quad $P^{p^j} \cdots P^pP\beta$
\note{exposants de la seconde ligne \uncertain{$P^{p^j}$} ; le début est en partie illisible.}

\[
  A_{\ast} = \mathrm{Sym}(\xi_i) \otimes \Lambda(\tau_j)
  \qquad \text{dual de } A
\]
\note{devant $A_{\ast}$, un mot biffé \struck{\ill{}} ; l'indice de $\tau$ est surchargé (\uncertain{$k$} corrigé en $j$).}
\[
  \Delta\xi_k = \sum_i \xi_{k-i}^{p^i}\otimes\xi_i
  \qquad
  \Delta\tau_k = \tau_k\otimes 1 + \sum_{i=0}^{k} \xi_{k-i}^{p^i}\otimes\tau_i
\]
p 85

\[
  A \otimes B \longrightarrow C \qquad \mathbb{F}_p\text{-modules}
\]
\[
  K(A,i) \wedge K(B,j) \longrightarrow K(C,i+j)
\]
cup-products pairings
\[
  K(\mathbb{F}_p,n)\wedge K(\mathbb{F}_p,m) \longrightarrow K(\mathbb{F}_p,m+n)
\]
\[
  H^{\ast}(K(\mathbb{F}_p,n))\otimes H^{\ast}(K(\mathbb{F}_p,m)) \longleftarrow
  H^{\ast}(K(\mathbb{F}_p,m+n),\mathbb{F}_p)
\]
\[
  A \otimes A \longleftarrow A
\]
\marginal{Le cup produit induit en coh la comultiplication sur l'alg de Steenrod $A$ correspondant à la multiplication dans la duale}

\[
  L\,\mathrm{Sym}^r = 0 \qquad r \neq p^s \qquad s \geq 0
\]
\[
  L\,\mathrm{Sym}^{p^s}K(\mathbb{F}_p,n)
\]
\[
  L_{n+i}\,\mathrm{Sym}\,K(\mathbb{F}_p,n) \xrightarrow[+1]{S}
  L_{n+i+1}\,\mathrm{Sym}\,K(\mathbb{F}_p,n+1)
\]
\[
  L^{\mathrm{st}}_{\ast}\mathrm{Sym}\,\mathbb{F}_p = A_{\ast}
\]
\[
  L^{\mathrm{st}}_{\ast}\,\text{\struck{$\mathbb{F}_p$}}\,[\mathbb{F}_p] =
\]
\note{les indices de $L$ dans la troisième ligne sont \uncertain{$n+i$} et \uncertain{$n+i+1$} ; le premier indice de la ligne $L^{\mathrm{st}}_{\ast}$ est surchargé.}

\uncertain{$1 \in$} $\mathrm{Sym}^{p^0}$ \qquad $\xi_1,\ \tau_0 \in \mathrm{Sym}^p$ \qquad $\xi$

\section{Construction bar}

\page{367}

\[
  A \xrightarrow{\ \varepsilon\ } \mathbb{F}_p
\]
$A$ $\mathbb{F}_p$-alg diff graduée aug sur $\mathbb{F}_p$
\[
  B(A) \qquad \overline{B}(A)
\]
\[
  A\otimes A\otimes A \rightrightarrows A\otimes A
  \underset{1\otimes\varepsilon}{\overset{\mu}{\rightrightarrows}} A
  \qquad \to \mathbb{Z}EG
\]
\[
  A\otimes A \rightrightarrows A \rightrightarrows \mathbb{F}_p
  \qquad \mathbb{Z}BG
\]
\note{les flèches multiples sont tracées à trois ou deux traits selon le degré ; une petite flèche courbe revient de $A$ vers $A\otimes A$ au-dessus de $\mu$.}

\marginal{$a\otimes\cdots\otimes a$, $n$ fois \quad (sous $A$ : $a$)}

$A$ commutative
\[
  \cdots\ \mathbb{Z}G^3 \rightrightarrows \mathbb{Z}G^2 \rightrightarrows \mathbb{Z}G
  \qquad \mathbb{Z}G^2 = \mathbb{Z}G\otimes\mathbb{Z}G
\]
Bar consts sur l'alg $\mathbb{Z}G$
\note{des flèches parallèles partent de $\mathbb{Z}G\otimes\mathbb{Z}G$ vers $\mathbb{Z}G$ et vers le bas à gauche ; le dessin n'est pas achevé.}

$A\otimes A$ \quad $A \to \mathbb{Z}$ \note{\uncertain{$\mathbb{Z}$} surchargé.}
\[
  A = \mathbb{Z}K(G,n) \qquad \overline{B}(A) \sim \mathbb{Z}BK(G,n)
\]
Structure d'algèbre de $\overline{B}(A)$

E Zilber
\[
  \overline{B}(A)\otimes\overline{B}(A) \rightleftarrows \overline{B}(A\otimes A)
  \longrightarrow \overline{B}(A)
\]
$A$ comm $\Longrightarrow$ $A\otimes A \to A$ hom d'alg
\note{devant $A\otimes A$, un signe \uncertain{$I\otimes$} ; « E Zilber » est sans doute Eilenberg-Zilber.}

\section{Le $K_2$ de Milnor}

\page{369}

Def de $K_2(R)$ de Milnor
\[
  E(R) \subset GL(R) \longrightarrow K_1(R) \to 0
\]
\emph{Th} \quad $E(R) = [E(R),E(R)]$
\note{le signe d'égalité est repassé ; \uncertain{$=$}.}

$R$ local \quad $E(R) = SL(R)$ \quad (Whitehead)
\[
  E(R) \hookrightarrow GL(R) \xrightarrow{\ \det\ } R^{\ast} \to 0
\]
\[
  0 \to K_2(R) \to St(R) \to E(R) \hookrightarrow Gl(R) \longrightarrow K_1(R) \to 0
\]
\note{une longue flèche va aussi directement de $St(R)$ à $Gl(R)$.}
\[
  \pi_1(E(R)) = K_2(R)
\]
\note{l'indice de $\pi$ est \uncertain{$1$}.}
\[
  BGl(R)^{+}
\]
\[
  \mathrm{Ext}^2(\mathbb{Z},G_m)
\]
\note{le premier argument de $\mathrm{Ext}^2$ est surchargé ; \uncertain{$\mathbb{Z}$}.}

Suite spectrale \qquad $X$ un schéma
\[
  H^i_{\mathrm{Zar}}(X,\underline{K}_j) \Longrightarrow K_{j-i}(X)
\]
\[
  H^i(X,\underline{K}_i) = i\text{-cycles}/\text{équiv rat}.
\]
\[
  0 \to G_m \to i_{\ast}G_{m,K} \to \bigoplus_{\text{pts de codim } 1}\mathbb{Z} \to 0
\]
\[
  0 \to \underline{K}_1 \to i_{\ast}\underline{K}_1 \to \bigoplus_{\text{pts } 1}\underline{K}_0 \to 0
\]
\[
  0 \to \underline{K}_i \to i_{\ast}\underline{K}_i \to \bigoplus_{\text{codim } 1} i_{x\ast}K_{i-1}
  \to \cdots \to \bigoplus_{\text{pts fermés}} i_{x\ast}\mathbb{Z} \to 0
\]
\[
  H^0(X,\ \cdots)\,/\,H^0(\quad)
\]
\note{au bas de la page, un quotient $H^0(X,\ldots)/H^0(\quad)$ dont le second argument du numérateur est \ill{} et le dénominateur laissé en blanc.}

\section{Cohomologie des espaces d'Eilenberg-MacLane}

\page{370}
\note{titre en haut à droite : « \uncertain{Coh} \struck{\ill{}} d'E-M I » ; la page suivante porte « Coh d'E-M II ».}

$H^{\ast}(K(\mathbb{F}_p,n),\mathbb{F}_p) = {}$ Algèbre anti com libre engendrée
par monômes admissibles d'excès $<n$ \qquad $p=2$ (alg sym)

\supplied{ou} « » \quad « » \quad $<n$ \quad $+$ « » \quad « » d'excès $= n$ qui se terminent par $\beta$ \qquad $p\neq 2$

\[
  A = \mathbb{F}_p[P^i,\beta]\,/\,P^0 = 1,\ \text{Adem}
  \qquad \deg P^i = 2i(p-1) \qquad \deg\beta = 1
\]
\[
  A = \mathbb{F}_2[\mathrm{Sq}^i]\,/\,\mathrm{Sq}^0 = 1,\ \text{Adem en car } 2
  \qquad \deg\mathrm{Sq}^i = i
\]
\[
  \mathrm{Sq}^a\mathrm{Sq}^b = \sum_{j=0}^{a<2b} (\quad)\,\mathrm{Sq}^{a+b-j}\mathrm{Sq}^j
\]
\note{les coefficients binomiaux de la relation d'Adem sont laissés en blanc.}

\marginal{$i>0$, $j\geq 0$ : $\deg\xi_i = 2p^i-2$, $\deg\tau_i = 2p^i-1$ ; $\mathrm{Sym}(\xi_i)\otimes\Lambda(\tau_j)$ $A_{\ast}$ l'algèbre duale. \quad $\mathrm{Sym}(\xi_i)$, $i>0$, $p=2$ : $\deg\xi_i = 2^i-1$.}

\begin{enumerate}
  \item $H_{\ast}(K(\mathbb{F}_p,n),\mathbb{F}_p)$
  \item $H_{\ast}(K(\mathbb{Z},n),\mathbb{F}_p)$
  \item $H_{\ast}(K(\mathbb{F}_p,n),\mathbb{Z})$
  \item $H_{\ast}(K(\mathbb{Z},n),\mathbb{Z})$
\end{enumerate}
\note{les numéros 1 et 2 sont cerclés sur la page.}

\[
  K(M\times N,n) \sim K(M,n)\times K(N,n)
\]

\marginal{\ill{} : $H_{\ast}(K(G,n),\mathbb{Z}) = \bigoplus G/pG \oplus {}_pG$, $G$ abélien quelconque. Ex : $H_{n+2}(K(G,n),\mathbb{Z}) = G/2G$ ; $H_{n+3}(K(G,n),\mathbb{Z}) = {}_2G$ ; $n+4$ : $= G/2G\oplus G/3G$ ; $n+5$ : $= {}_2G + {}_3G$, dans la partie stable $H_{n+i}(K(G,n))$, $0\leq i<n$. \quad \uncertain{Deprime}}

\[
  G \in \mathcal{R} \text{ des car } p \qquad
  H^{\ast}(K(G,n),\mathcal{R})
\]
alg anticom libre sur le module $\mathrm{Hom}(G,\mathcal{R})\otimes_{\mathbb{F}_p}$ (partie de $A$ d'excès $<n$)
\note{« une $\mathbb{F}_p$-algèbre » est ajouté au-dessus de $\mathcal{R}$, au-dessus du mot biffé \struck{groupes} ; après « module » un mot biffé \struck{\ill{}}.}

l'homologie est entièrement de \emph{torsion première}.

Bockstein
\[
  \mathbb{Z} \xrightarrow{\ p\ } \mathbb{Z} \to \mathbb{Z}/p
\]
\[
  \xrightarrow{0 = p} H_{\ast}(K(\mathbb{F}_p,n),\mathbb{Z}) \to H_{\ast}(K(\mathbb{F}_p),\mathbb{F}_p)
  \to H_{\ast-1}(K(\mathbb{F}_p),\mathbb{Z}) \xrightarrow{\ p=0\ } 0
\]
\note{une flèche oblique étiquetée $\beta_1$ va de $H_{\ast}(K(\mathbb{F}_p),\mathbb{F}_p)$ vers $H_{\ast-1}(K(\mathbb{F}_p),\mathbb{F}_p)$, qui est écrit sous $H_{\ast-1}(K(\mathbb{F}_p),\mathbb{Z})$ et relié à lui par une flèche verticale.}

Puisque mult par $p = 0$, il suffit de \struck{\ill{}} expliciter $\beta_1$ sur l'algèbre $A_{\ast}$
\[
  K(\mathbb{Z},n) \to K(\mathbb{Z},n) \to K(\mathbb{F}_p,n)
\]
induit en deg $<2n$ une suite exacte (à partir de suite de Leray) qui est exacte pour la même raison que ci-dessus :
\[
  0 \to H_{\ast}(K(\mathbb{Z},n),\mathbb{F}_p) \to H_{\ast}(K(\mathbb{F}_p,n),\mathbb{F}_p)
  \to H_{\ast-1}(K(\mathbb{Z},n),\mathbb{F}_p) \to 0
\]
\note{une flèche oblique étiquetée $\beta_2$ va de $H_{\ast}(K(\mathbb{F}_p,n),\mathbb{F}_p)$ vers $H_{\ast-1}(K(\mathbb{F}_p,n),\mathbb{F}_p)$, écrit sous le troisième terme.}

les 2 $\beta_i$ \struck{est une} \add{sont des} dérivation sur l'algèbre graduée $A_{\ast}$
\[
  \mathrm{Sym}(\xi_i)\otimes\Lambda(\tau_i)
\]
l'un des 2 $\beta_i$ :
\[
  \tau_0 \to 1 \qquad
  \begin{cases} \tau_i \to 0 & i\neq 0 \\ \xi_j \to 0 & j>0 \end{cases}
\]

Conclusion \quad $H_{\ast}(K(\mathbb{F}_p,n),\mathbb{Z}) = {}$ \struck{l'autre} l'alg à puiss div libre eng par la partie d'excès $<n$ de $\{\mathrm{Sym}(\xi_i)\otimes\Lambda(\tau_j)\}$, $i>0$, $j\geq 0$.
\note{dans cette ligne et la suivante, le coefficient $\mathbb{Z}$ est repassé ou biffé ; lecture \uncertain{$\mathbb{Z}$}.}

$H_{\ast}(K(\mathbb{Z},n),\mathbb{Z})$ \quad $\beta_2$ :
\[
  \begin{cases} \tau_j \to \xi_j \\ \xi_j \to 0 \end{cases}
\]

$\beta_2$ agit sur $H_{\ast}(K(\mathbb{F}_p,n),\mathbb{Z}) \xrightarrow{\ \beta_2\ } H_{\ast}(K(\mathbb{F}_p,n),\mathbb{Z})$.
C'est une dérivation (anti com) :
\[
  \tau_j \longmapsto \xi_j \qquad \xi_j \longmapsto 0
\]
\[
  \mathrm{Im}\,\beta_2 = \ker\beta_2 = H_{\ast}(K(\mathbb{Z},n),\mathbb{Z})_{(p)} .
\]

La conjugaison $c : A \to A$, $: A_{\ast} \to A_{\ast}$ est explicitée dans Steenrod. Elle transforme l'un des $\beta$ dans l'autre et transporte la sous algèbre (partie stable de) :
\[
  H_{\ast}(K(\mathbb{Z},n),\mathbb{F}_p) \xrightarrow{\ \sim\ } H_{\ast}(K(\mathbb{F}_p,n),\mathbb{Z})
\]
\[
  (\ast) \qquad H_{\ast}(K(\mathbb{F}_p,n),\mathbb{F}_p) = A_{\ast} \longleftarrow
\]
\note{$H_{\ast}(K(\mathbb{Z},n),\mathbb{F}_p)$ est relié par une inclusion à $A_{\ast}$ ; une flèche courbe revient de la droite vers $A_{\ast}$. Le repère « $(\ast)$ » est repris p. 371.}

\page{371}
\note{titre en haut à droite : « Coh d'EM II ».}

\emph{Th} (Cartan) : $G$, $H$ deux \struck{anneaux} \add{groupes} \struck{\ill{} $\mathbb{F}_p$-modules} (il fait des hyp sur $G$ et $H$).
\[
  H_{\ast}(K(G,n),H) \simeq H_{\ast}(K(H,n),G) \quad \text{additif.}
\]
\emph{Dém} : Calcul.

$X$ un ensemble simplicial pointé.
\[
  \pi_i(\mathbb{Z}X) \simeq H_i(X)
\]
\[
  K(\mathbb{Z},n)\wedge X \xrightarrow{\ A(X)\ } \mathbb{Z}^{+}(S^nX)
  \qquad \mathbb{Z}^{+}(X,x) = \mathbb{Z}X/\mathbb{Z}x
\]
\note{devant $\mathbb{Z}^{+}(S^nX)$, un mot biffé \struck{\ill{}}.}
$2n-1$ quasi isomorphisme.
\[
  K(\mathbb{Z},n) = \mathbb{Z}^{+}[S^n]
\]
\[
  \left.\begin{aligned}
    \mathbb{Z}(X\times Y) &\simeq \mathbb{Z}X\otimes\mathbb{Z}Y \\
    \mathbb{Z}^{+}(X\wedge Y) &\simeq \mathbb{Z}^{+}X\otimes\mathbb{Z}^{+}Y
  \end{aligned}\right\}
  \ \begin{array}{l} \text{EZ} \\ \text{EZ pointé.} \end{array}
\]
\[
  K(\mathbb{Z},n)\wedge X \to K(\mathbb{Z},n)\wedge\mathbb{Z}^{+}X
\]
\[
  \mathbb{Z}^{+}(S^n)\wedge X \to \mathbb{Z}^{+}(S^n)\wedge\mathbb{Z}^{+}X
  \to \mathbb{Z}^{+}(S^n)\otimes\mathbb{Z}^{+}X \xrightarrow{\ \mathrm{EZ}\ } \mathbb{Z}^{+}(S^nX)
\]
\[
  \pi_{n+r}(K(\mathbb{Z},n)\wedge X) = \pi_{n+r}(\mathbb{Z}^{+}(S^nX))
  = \widetilde{H}_{n+r}(S^nX) = \widetilde{H}_r(X)
\]
\emph{Dém} : $A(S^n)$
\[
  \pi_i(K(\mathbb{Z},n)\wedge S^m) \qquad
  \mathbb{Z}^{+}(S^n\wedge S^m) = \mathbb{Z}^{+}(S^{n+m})
\]
\[
  S^m\wedge K(\mathbb{Z},n) \xrightarrow{\ \mathrm{qis}\ } K(\mathbb{Z},n+m)
\]
\[
  S^m \to \mathbb{Z}^{+}(S^m) = K(\mathbb{Z},m)
\]
\[
  S^m\wedge K(\mathbb{Z},n) \to K(\mathbb{Z},m)\wedge K(\mathbb{Z},n)
  \xrightarrow{\ \cup\ } K(\mathbb{Z},n+m).
\]
\uncertain{Pt clef} : \quad $X \xrightarrow{\ h_r(\ )\ } \pi_{n+r}(K(\mathbb{Z},n)\wedge X)$ est une th d'homologie partielle
\[
  Y \hookrightarrow X \to X/Y
\]
induit
\[
  h_r(Y) \to h_r(X) \to h_r(X/Y) \to h_{r-1}(Y) \to \cdots
\]
\emph{$r$ assez petit.}

\note{encadré, en bas à gauche :}
Ex de la dém 2 : interprétation du diagramme $(\ast)$ en bas de la page I.
\[
  \begin{array}{ccc}
    K(\mathbb{Z},n)\wedge K(\mathbb{F}_p,m) & \longrightarrow & K(\mathbb{F}_p,m)\wedge K(\mathbb{Z},n) \\
    \big\downarrow & & \big\downarrow \\
    K(\mathbb{F}_p,n)\wedge K(\mathbb{F}_p,m) & \longrightarrow & K(\mathbb{F}_p,n)\wedge K(\mathbb{F}_p,m)
  \end{array}
\]
\note{dans l'encadré, les indices $n$, $m$ sont surchargés à la ligne du bas, et la flèche verticale de gauche est un crochet d'inclusion ; « Pt clef » est écrit au bord de l'encadré.}

Dém 2 du th de Cartan
\[
  \pi_{i+m}(K(H,m)\wedge K(G,n)) \simeq \pi_{i+n}(K(G,n)\wedge K(H,m))
\]
s'obtient par permutation des facteurs
\note{le second terme est écrit \uncertain{$\pi_{i+n}(K(G,n)\wedge K(H,n))$}.}

\section{Résolutions canoniques d'un groupe abélien}

\page{372}

\[
  X_{\ast}(G) \longrightarrow G \qquad X_n(G) = \mathbb{Z}(G^n\times\mathbb{Z}^{+})
\]
\note{lecture de $X_n(G)$ \uncertain{incertaine}.}
\[
  \mathbb{Z}^{\mathbb{Z}^G} \rightrightarrows \mathbb{Z}^G \qquad G
  \qquad\qquad G \to G\,! \qquad H
\]
\[
  \mathbb{Z}^{\mathbb{Z}K(G,n)} \rightrightarrows \mathbb{Z}^{K(G,n)}
  \qquad K(G,n) = G\otimes_{\mathbb{Z}}\Lambda^n\Psi
\]
\marginal{Pour tout $n$ on a une résolution simpliciale de $K(G,n)$.}
$n$ \struck{large} grand

$r$-qis \quad $r < f(n,p)$
\[
  K(\mathbb{Z},n)\wedge K(G,p) \xrightarrow{\ \sim\ } \mathbb{Z}(S^n\wedge K(G,p))
  \longrightarrow \mathbb{Z}K(G,n+p)
\]
\[
  S^n\wedge K(G,p) \to K(G,n+p)
\]
\[
  K(\mathbb{Z},p_1)\wedge K(\mathbb{Z},p_2)\wedge\cdots\wedge K(\mathbb{Z},p_r)\wedge K(G,p_{r+1})
  \to \mathbb{Z}^{r}K(G,\ p)
\]
\[
  \sum_{i=1}^{r+1} p_i = p
\]
\note{dans $\mathbb{Z}^rK(G,p)$, une lettre biffée devant $p$.}

Encadré :
\[
  r\text{ qis} \qquad K(\mathbb{Z},n)\wedge K(G,p) \longrightarrow \mathbb{Z}K(G,n+p)
\]
\[
  r = \min(n+2p-1,\ 2n+p-1)
\]
\[
  K(\mathbb{Z},n)\wedge X \to \mathbb{Z}[S^n\wedge X]
\]
\struck{\ill{}}
\[
  K(\mathbb{Z},n)\wedge K(\mathbb{Z},m)\wedge K(G,p) \to \mathbb{Z}[S^n\wedge K(\mathbb{Z},m)\wedge K(G,p)]
\]
\[
  \downarrow
\]
\[
  \mathbb{Z}[K(\mathbb{Z},n+m)\wedge K(G,p)]
\]
\[
  \downarrow
\]
\[
  \mathbb{Z}[\mathbb{Z}[S^{n+m}\wedge K(G,p)]]
\]
\[
  \downarrow
\]
\[
  \mathbb{Z}[\mathbb{Z}K(G,n+m+p)]
\]
\note{dans la troisième ligne de la colonne, une lettre biffée après le second $\mathbb{Z}$ ; le $G$ de $K(G,p)$ en tête de ligne surcharge un $\mathbb{Z}$.}

\section{Foncteurs dérivés de Sym}

\page{373}
\note{titre en haut à droite, souligné deux fois : « Foncteurs dérivés de Sym ». La lettre que nous rendons par $\mathcal{O}$ dans la seconde moitié de la page est un O cerclé ou barré, écrite en surcharge sur $G_a$ ; lecture \uncertain{$\mathcal{O}$}.}

\[
  H^q(X,R) \simeq H^q(X,\mathbb{F}_p)\otimes_{\mathbb{F}_p} R
\]
$\Theta\in\mathcal{A} = {}$ algèbre de Steenrod
\[
  P^i_R = P^i\otimes 1 \qquad i>0 \qquad 2i(p-1)
\]
\[
  P^0_R = P^0\otimes F \qquad (P^0 = 1)
\]
\marginal{sous $P^i_R$ : construit par Steenrod.}
\note{$F$ est sans doute le Frobenius ; $P^0_R$ est cerclé au crayon.}

$G_a = {}$ Fais structural de $(\mathrm{Sch}/\mathbb{F}_p)$ \qquad $G_a\otimes(\Lambda^n\Psi)$
\[
  H^{\ast}(K(G_a,n),G_a) = \mathcal{A}\cdot\text{les op}
\]
\[
  \mathbb{H}^{\ast}(\mathbb{Z}[K(G_a,n)]^{\sim},G_a)
\]
Illusie \quad $X\in\mathrm{ob}\,D_\bullet(T)$ \quad cat dérivée des fais simpl \quad $F$ objet abélien
\[
  \mathbb{H}^{\ast}(X,F) = [X,K(F,n)]
  \qquad X' \to K(F,n),\ X' \to X
\]
\note{une flèche monte de $H^j(X_i,F)$ vers $\mathbb{H}^{\ast}(X,F)$ ; le croquis hachuré en haut à droite de la page n'est pas reproduit.}

Epstein \uncertain{$\mathcal{L}$}. Operations in homological algebra, Invent. Math. 1 (1966)

Illusie p 69 \quad 3.2.1.19 -- 3.2.1.18
\[
  \mathcal{E}xt^p(\Lambda(X),F) = \varinjlim_{\substack{X'\to X\\ \text{qis}\\ \text{fibrant}}} [X',K(F,p)]
\]
\[
  \overset{\text{déf}}{=} \mathbb{H}^p(X,F)
\]
\[
  \mathcal{E}xt^p(\Lambda(-),F) \to \mathrm{Ext}^q(\Lambda(\ ),G)
\]
Op coh de type $(F,p,G,q)$ sont class par $\mathbb{H}^q(K(F,p),G)$
\[
  \mathbb{Z}X \xrightarrow{\ f\ } K(A,q)
\]

\[
  \mathbb{H}^n(K(G_a,m),G_a) = R^n\,\mathrm{Sym}(\mathbb{F}_p[m])
  = \text{dual de } L^n\Gamma\,(\mathbb{F}_p[m])
\]
\note{devant $(\mathbb{F}_p[m])$, un mot biffé \struck{Sym} ; la lecture $L^n\Gamma$ est \uncertain{incertaine}.}

\emph{Théorème} : Soit $\mathcal{A}$ l'algèbre sur $\mathbb{F}_p$ eng par $P^i$, $\deg P^i = 2i(p-1)$, $i\geq 0$, sat aux relations d'Adem (mais non à $P^0 = 1$).

$\mathbb{H}^{\ast}(K(\mathcal{O},m),\mathcal{O})$ est l'alg anticomm libre engendrée par les monômes admissibles de $\mathcal{A}$ d'excès $<n$.

\emph{Dém} : Comme dans le cas classique, il suffit de calculer $H^{\ast}(K(G_a,1),G_a)$ et de grimper en utilisant le th de Borel et la suite sp de Leray associée à
\[
  K(G_a,n-1) \to \ast \to K(G_a,n)
\]
\note{$\ast$ est écrit au-dessus de $K(G_a,n)$, avec une flèche verticale.}

Lazard : Coh de l'analyseur classique dans « Lois de groupes formels et analyseurs » Ann ENS
\note{une flèche relie $H^{\ast}(K(G_a,1),G_a)$ à la référence de Lazard.}

\[
  G_a^3 \qquad G_a\times G_a \qquad G_a \qquad\qquad
  k[x,y,z] \leftleftarrows k[x,y] \leftleftarrows k[x]
\]
\note{au-dessus de $k[x]$ : \uncertain{$\mathrm{Sym}\,\mathbb{F}_p$} ; les flèches sont à plusieurs traits.}
\[
  \iota_1 \in H^1(K(G_a,1),G_a)
\]
\[
  \text{Witt} \leftarrow \beta \in H^2(K(G_a,1),G_a)
\]
\[
  0 \to G_a \to W_2 \to G_a \to 0 \qquad
  0 \to \mathbb{Z}/p \to \mathbb{Z}/p^2 \to \mathbb{Z}/p \to 0
\]

\[
  \Bigl|\ H^{\ast}(K(\mathbb{F}_{p^h},m),\mathbb{F}_{p^h})
\]
$R$ anneau de car $p\neq 0$
\[
  H^{\ast}(K(R,n),R) = {}
\]
Alg anticomm libre $/R$ eng par $\mathrm{Hom}(R,R)$-module libre de eng par les monômes admissibles
\[
  \bigoplus P^If\ \ i_n \qquad f\in\mathrm{Hom}(R,R)
\]
\[
  H^{\ast}(X,\mathcal{O}) \to H^{\ast}(X,\mathcal{O}) \qquad x \mapsto P^Ix
\]
$\mathrm{Hom}$

On se ramène au calcul de $H^{\ast}(K(\mathbb{F}_{p^h},1),\mathbb{F}_{p^h})$
\[
  x\in H^{\ast}(K(\mathcal{O},m),\mathcal{O})
  \qquad K(\mathcal{O},m) \longrightarrow K(\mathcal{O},n)
  \qquad \forall\, R/\mathbb{F}_p
\]
\[
  K(R,m) \longrightarrow K(R,n)
\]
\[
  x\in H^{\ast}(K(\mathcal{O},n),\mathcal{O})_h \xrightarrow{\ \text{injective}\ }
  \mathbb{H}^{\ast}(K(\mathbb{F}_{p^h},n),\mathbb{F}_{p^h})
\]
représenté par des polynômes de deg $<p^h$ en chaque var

\section{Calcul des $\mathrm{Ext}^i(G_a,G_a)$}

\page{374}

\emph{Prop} : Le $k[F]$-module $\mathrm{Ext}^i(G_a,G_a)$ (pour l'action de $\mathrm{Hom}(G_a,G_a)$ à droite) est de $F$ torsion ($i>1$).
\note{la borne $i>1$ surcharge un autre chiffre.}

\emph{Cor} : (\uncertain{Milne}) \quad $G_a^{\mathrm{parf}} = {}$ restriction de $G_a$ au site parfait$/k$
\[
  G_a^{\mathrm{parf}} = \varinjlim_{F} G_a
\]
\[
  \mathrm{Ext}^i(G_a^{\mathrm{parf}},G_a^{\mathrm{parf}}) = 0 \qquad i>1.
\]
\note{la flèche sous « lim » a deux pointes ; nous la lisons comme une limite inductive. Une longue flèche courbe relie la Prop à la Dém.}

\emph{Dém} : On considère la suite spectrale du coefficient universel
\[
  \forall \text{ fais simpl } X : \quad
  E_2^{i,j} = \mathrm{Ext}^i_{\mathbb{F}_p}(\underline{H}_j(X),F) \Longrightarrow \mathbb{H}^{i+j}(X,F)
\]
En fait on va travailler dans la catégorie des fais de $\mathbb{F}_p$-modules

On travaille mod $F$ (au sens mod $\mathcal{C}$)
\[
  \mathrm{Hom}(G_a,G_a) = k[F] \qquad \mathrm{Ext}^1(G_a,G_a) = 0
\]
\note{à droite, biffé : \struck{$E^{\ast i} = 0$, $j<n$}.}

Dém par récurrence : On va considérer la s. spectrale précédente avec $X = K(G_a,n)$ \qquad \uncertain{$\mathrm{Ext}^i$} \quad $n$ grand.
\[
  E \qquad \mathrm{Ext}^i(H_j(K(G_a,n)),G_a) \Longrightarrow \mathbb{H}^{\ast}(K(G_a,n),G_a)
\]
\[
  E_2^{i,n} = \mathrm{Ext}^i(H_n(K(G_a,n)),G_a) = \mathrm{Ext}^i(G_a,G_a)
\]
\drawing{à droite, un petit diagramme du plan $(i,j)$ : la ligne $j=n$ porte des $0$, et la région au-dessus est hachurée.}
\[
  \underline{H}_j(K(G_a,n)) = \bigoplus G_a \qquad G_a \xrightarrow{\ p=0\ } G_a
\]

\page{375}
\note{feuille à portées, écrite en largeur ; suite directe de la p. 374.}

\[
  \mathrm{Ext}^i(H_j(K(G_a,n)),G_a) \Longrightarrow \mathbb{H}^{i+j}(K(G_a,n),G_a)
\]
On raisonne par récurrence, mod $F$
\[
  \mathrm{Ext}^i(H_n(K(G_a,n)),G_a) \to H^{n+i}(K(G_a,n),G_a)
  \to \mathrm{Hom}(H_{n+i}(K(G_a,n)),G_a)
\]
\[
  \to \mathrm{Ext}^{i+1}(G_a,G_a) \xrightarrow{\ \partial\ } H^{n+i+1}(K(G_a,n),G_a)
\]
\note{sous $\mathrm{Hom}(H_{n+i}(K(G_a,n)),G_a)$ : « $= \bigoplus\mathrm{Hom}(G_a,G_a)$ » ; un $0$ avec une flèche oblique vers $H^{n+i}$. Sous la seconde ligne, une flèche verticale descend de $\mathrm{Ext}^{i+1}(G_a,G_a)$ et une flèche montante relie les deux groupes de droite :}
\[
  \mathrm{Ext}^{i+1}(\mathbb{F}_{p^h},\mathbb{F}_{p^h}) \xrightarrow{\ \partial\ }
  H^{n+i+1}(K(\mathbb{F}_{p^h},n),\mathbb{F}_{p^h})
  \qquad \mathrm{Ext}^{i+1}(\mathbb{F}_{p^h},\mathbb{F}_{p^h}) = 0
\]

\emph{Lemme 1} \quad le edge homomorphism $\partial$ est nul.

\emph{Lemme 2} \quad $x$ un générateur de $H^{n+i}(K(G_a,n),G_a) \to \mathrm{Hom}(H_{n+i}(K(G_a,n)),G_a)$
\[
  \simeq \mathrm{Hom}(G_a\otimes H_{n+i}(K(\mathbb{F}_p,n)),G_a)
  = H^{n+i}(K(\mathbb{F}_p,n),\mathbb{F}_p)\otimes\mathrm{Hom}(G_a,G_a)
\]
\[
  P^I\iota_n \qquad \beta P^{i_1}\beta\cdots P^{i_k}\iota_n \qquad \text{longueur } n+1
\]
\[
  P^d\iota_n \longrightarrow P^d\iota_n\otimes F
\]
\note{les exposants de $\beta P^{i_1}\beta\cdots P^{i_k}$ sont \uncertain{incertains}, et une accolade souligne une partie du monôme.}

Dém \quad \uncertain{on} le démontre pour $P^0\iota_n$

\[
  \mathrm{Ext}^{\ast}_{\mathbb{F}_p}(G_a,G_a) \qquad
  \mathrm{Ext}^{2i+1}(G_a,G_a) = 0 \quad \forall i \quad i>0
\]
\[
  \mathrm{Ext}^{2i}(G_a,G_a) = k[F]/(F^d) \qquad d = v_p(i).
\]
\[
  \mathcal{A} \to A \qquad P^i \mapsto P^i \qquad P^0 \mapsto 1
\]
\[
  \mathcal{A} \to \mathcal{A}\otimes\mathcal{A} \qquad P^i \mapsto \sum_{j+k=i} P^j\otimes P^k
\]
\[
  \mathbb{Z}^iK(G_a,n) \cdots \mathbb{Z}\mathbb{Z}K(G_a,n) \rightrightarrows \mathbb{Z}K(G_a,n)
\]
\note{les flèches multiples du complexe sont tracées à deux ou trois traits.}
\[
  E_1 \qquad \cdots\ \mathcal{A}^{\otimes 2}\otimes\mathcal{A} \leftleftarrows \mathcal{A}\otimes\mathcal{A} \leftleftarrows \mathcal{A}
\]
Dualement
\[
  \cdots\ \mathcal{A}_{\ast}^{\otimes 2}\otimes\mathcal{A}_{\ast} \rightrightarrows \mathcal{A}_{\ast}\otimes\mathcal{A}_{\ast} \rightrightarrows \mathcal{A}_{\ast}
\]
\[
  E^2 = \mathrm{Tor}^{A_{\ast}}(\mathcal{A}_{\ast},\mathbb{F}_p) \Longrightarrow E^2 \text{ par Koszul}
\]
\note{la ligne $E^2 = \mathrm{Tor}\ldots$ est \uncertain{incertaine} à gauche (« $E^2$ » ou « $E^1$ »).}

\[
  \mathrm{Ext}^i(\mathbb{Z}^jK(G_a,n),G_a) \Longrightarrow \mathrm{Ext}^{i+j}(G_a,G_a)
\]
$K(\mathbb{Z}) = K(\mathbb{Z},m)$ \quad $n$ grand
\[
  \mathbb{Z}[\mathbb{Z}^{j-1}K(G_a,n)] \sim \mathbb{Z}[K(\mathbb{Z})\wedge K(\mathbb{Z})\wedge\cdots\wedge K(G_a,n)]
\]
\[
  E_2^{i,j} = \widetilde{\mathbb{H}}^{i}(K(\mathbb{Z})^{\wedge j-1}\wedge K(G_a),G_a)
  \overset{\text{Künneth}}{\sim}
  \widetilde{\mathbb{H}}^{\ast}(K(\mathbb{Z}))^{\otimes j-1}\otimes\widetilde{\mathbb{H}}^{\ast}(K(G_a),G_a)
\]
\[
  \mathbb{H}^{\ast}(K(\mathbb{F}_p),\mathbb{F}_p) = A \qquad
  \mathbb{H}^{\ast}(K(G_a),G_a) = \mathcal{A}
\]
\note{l'argument se poursuit au-delà de cette page, hors du lot.}

\end{document}
